\chapter{Verwandte industrielle Systeme}

\section{Einordnung industrieller Retrieval- und Ranking-Systeme}

Die im vorliegenden Projekt eingesetzte Architektur basiert auf etablierten Verfahren des Information Retrieval und Learning-to-Rank.
Vergleichbare Architekturen finden sich seit vielen Jahren in Suchmaschinen, Enterprise-Search-Plattformen, Produktsuchen sowie Empfehlungssystemen.
Obwohl moderne Systeme zunehmend neuronale Verfahren und Vektorsuchen integrieren, bilden die in Kapitel \ref{chap:stand_der_technik} beschriebenen Komponenten weiterhin einen wesentlichen Bestandteil produktiver Suchsysteme.
Insbesondere die Kombination aus

\begin{itemize}
  \item Kandidatenretrieval,
  \item Rank Fusion,
  \item Feature Engineering,
  \item Learning-to-Rank,
  \item Re-Ranking sowie
  \item Ergebnisaggregation
\end{itemize}

stellt bis heute die dominierende Architektur vieler industrieller Retrievalsysteme dar.

\section{Microsoft Bing}

Microsoft Bing gilt als eines der bekanntesten produktiven Learning-to-Rank-Systeme und spielte eine zentrale Rolle bei der Entwicklung der RankNet-, LambdaRank- und LambdaMART-Verfahren \cite{burges2010lambdamart}.
Die Bing-Sucharchitektur besteht aus mehreren aufeinanderfolgenden Retrieval- und Rankingstufen:

\begin{enumerate}
  \item Candidate Retrieval
  \item Feature Extraction
  \item Learning-to-Rank
  \item Re-Ranking
  \item Ergebnispräsentation
\end{enumerate}

Die initiale Kandidatenselektion erfolgt mithilfe effizienter Retrievalverfahren.
Die erzeugten Kandidaten werden anschließend durch eine Vielzahl strukturierter Merkmale beschrieben.
Ein Learning-to-Rank-Modell berechnet daraus eine Relevanzbewertung.
Die grundlegende Architektur entspricht damit dem folgenden Aufbau:

\begin{align}
  Query \rightarrow Retrieval \rightarrow Feature\ Extraction \rightarrow LTR \rightarrow Ranking
\end{align}

Dieses Architekturprinzip bildet die Grundlage zahlreicher moderner Suchmaschinen.

\section{Azure AI Search}

Azure AI Search stellt die Enterprise-Search-Plattform von Microsoft dar.
Das System kombiniert mehrere Retrievalverfahren innerhalb einer Hybrid-Search-Architektur.
Dabei werden unterschiedliche Suchverfahren parallel ausgeführt und anschließend durch Rank-Fusion zusammengeführt.
Die Architektur umfasst typischerweise:

\begin{itemize}
  \item Volltextsuche auf Basis von BM25,
  \item semantische Suchverfahren,
  \item Hybrid Search,
  \item Rank Fusion mittels Reciprocal Rank Fusion,
  \item zusätzliches Re-Ranking.
\end{itemize}

Die verwendete Rank-Fusion basiert auf demselben RRF-Verfahren, das durch Cormack et al. beschrieben wurde \cite{cormack2009rrf}.
Im Gegensatz zur klassischen Vektorsuche besitzt dieser Ansatz den Vorteil, dass sowohl exakte Merkmalstreffer als auch ähnliche Dokumente berücksichtigt werden können.
Anwendungsbereiche umfassen:

\begin{itemize}
  \item technische Dokumentensuche,
  \item Wissensdatenbanken,
  \item Produktkataloge,
  \item Serviceportale,
  \item unternehmensweite Dokumentensuche.
\end{itemize}

\section{Elasticsearch Learning-to-Rank}

Elasticsearch zählt zu den am weitesten verbreiteten Suchplattformen im industriellen Umfeld.
Die optionale Learning-to-Rank-Erweiterung implementiert nahezu dieselbe Architektur wie das vorliegende Projekt:

\begin{enumerate}
  \item BM25-Kandidatenretrieval
  \item Merkmalsextraktion
  \item Learning-to-Rank-Modell
  \item Re-Ranking der Kandidatenliste
\end{enumerate}

Das zugrunde liegende Modell wird üblicherweise außerhalb von Elasticsearch trainiert und anschließend als Rankingmodell eingebunden.
Typische Merkmale umfassen:

\begin{itemize}
  \item Termstatistiken,
  \item BM25-Scores,
  \item Dokumentattribute,
  \item Kategoriesignale,
  \item numerische Ähnlichkeitsmaße.
\end{itemize}

Die resultierende Architektur entspricht praktisch dem klassischen LETOR-Ansatz nach Liu \cite{liu2009learning}.
Einsatzgebiete finden sich insbesondere in

\begin{itemize}
  \item E-Commerce-Plattformen,
  \item Produktkatalogen,
  \item Dokumentensuchen,
  \item technischen Wissensdatenbanken.
\end{itemize}

\section{Apache Solr Learning-to-Rank}

Apache Solr implementiert ein Learning-to-Rank-Framework, welches ähnlich wie Elasticsearch aufgebaut ist.
Die Pipeline besteht aus:

\begin{enumerate}
  \item Dokumentretrieval,
  \item Featuregenerierung,
  \item Rankingmodell,
  \item Re-Ranking.
\end{enumerate}

Unterstützt werden unter anderem:

\begin{itemize}
  \item Gradient Boosted Trees,
  \item LambdaMART,
  \item externe Machine-Learning-Modelle.
\end{itemize}

Solr wird häufig eingesetzt für:

\begin{itemize}
  \item E-Commerce,
  \item Produktfindung,
  \item Unternehmenssuche,
  \item Wissensmanagementsysteme.
\end{itemize}

Die technische Nähe zum vorliegenden Projekt ist insbesondere bei der Nutzung expliziter strukturierter Merkmale sehr hoch.

\section{Amazon Produktsuche}

Die Produktsuche von Amazon zählt zu den bekanntesten kommerziellen Suchsystemen.
Aufgrund der Größe des Produktkataloges müssen Suchanfragen in mehreren Stufen verarbeitet werden.
Typischerweise erfolgt:

\begin{enumerate}
  \item schnelles Kandidatenretrieval,
  \item Ranking über Produktmerkmale,
  \item Personalisierung,
  \item Re-Ranking.
\end{enumerate}

Für Produktsuchen besitzen strukturierte Merkmale eine besondere Bedeutung:

\begin{itemize}
  \item Kategorie,
  \item Hersteller,
  \item technische Eigenschaften,
  \item Preise,
  \item Verfügbarkeit,
  \item Produktspezifikationen.
\end{itemize}

Somit ähneln Produktsuchen dem hier betrachteten technischen Artikel-Makro-Matching wesentlich stärker als allgemeine Websuchmaschinen.
Insbesondere das intensive Feature Engineering und die Nutzung Gradient-Boosted-Tree-basierter Rankingverfahren gelten in Produktsuchen seit vielen Jahren als etablierter Standard.

\section{eBay Cassini}

Cassini ist die Suchplattform von eBay.
Das System nutzt eine mehrstufige Rankingpipeline zur Bewertung von Produktangeboten.
Berücksichtigt werden unter anderem:

\begin{itemize}
  \item Suchbegriffe,
  \item Kategorien,
  \item Angebotsattribute,
  \item Verkaufsmetriken,
  \item Qualitätsindikatoren.
\end{itemize}

Die Architektur folgt ebenfalls dem Retrieval- und Re-Ranking-Paradigma und nutzt zahlreiche strukturierte Merkmale zur Bestimmung der Ergebnisreihenfolge.
Damit weist Cassini starke Parallelen zu klassischen Learning-to-Rank-Systemen auf.

\section{Vergleich der Systeme}

Die Tabelle zeigt die Ähnlichkeit zwischen den betrachteten industriellen Systemen und der entwickelten Architektur.

\begin{table}[H]
  \centering
  \caption{Vergleich industrieller Suchsysteme}
  \begin{tabular}{lccccc}
    \toprule
    System & BM25 & RRF & LTR & Re-Ranking & Aggregation \\
    \midrule
    Bing & \checkmark & - & \checkmark & \checkmark & \checkmark \\
    Azure AI Search & \checkmark & \checkmark & \checkmark & \checkmark & - \\
    Elasticsearch LTR & \checkmark & - & \checkmark & \checkmark & - \\
    Apache Solr LTR & \checkmark & - & \checkmark & \checkmark & - \\
    Amazon Search & \checkmark & - & \checkmark & \checkmark & - \\
    eBay Cassini & \checkmark & - & \checkmark & \checkmark & - \\
    Vorliegendes System & \checkmark & \checkmark & \checkmark & \checkmark & \checkmark \\
    \bottomrule
  \end{tabular}
\end{table}

\section{Zusammenfassung}

Die Analyse verwandter industrieller Systeme zeigt, dass die entwickelte Architektur keinen experimentellen Ansatz darstellt, sondern auf etablierten Retrieval- und Rankingverfahren basiert.
Insbesondere die Kombination aus

\begin{itemize}
  \item BM25-Retrieval,
  \item Rank Fusion mittels RRF,
  \item strukturierter Merkmalsextraktion,
  \item LambdaMART-basiertem Learning-to-Rank,
  \item mehrstufigem Re-Ranking sowie
  \item Evidenzaggregation
\end{itemize}

entspricht einem Architekturmuster, das seit vielen Jahren in Suchmaschinen, Produktsuchen und Enterprise-Search-Systemen erfolgreich eingesetzt wird.
Die entwickelte Lösung kann daher als fachlich fundierte, industriell etablierte und wissenschaftlich abgesicherte Retrieval- und Rankingarchitektur eingeordnet werden.
